\chapter{اللغة التقاربية والجمع الجزئي لأبيل}

\section{لماذا نحتاج إلى لغة تقاربية؟}

تهتم نظرية الأعداد التحليلية غالبًا بسلوك كميات حسابية عندما تكبر المعلمة الأساسية. وقد يكون إيجاد صيغة مغلقة دقيقة أمرًا مستحيلًا أو غير مفيد، بينما يمكن فصل الكمية إلى حد رئيسي يصف السلوك المنتظم وحد خطأ يقيس التذبذب.

إذا كانت \(F(x)\) دالة عدّ أو مجموعًا حسابيًا، فإن الصيغة النموذجية تكون
\[
F(x)=M(x)+E(x),
\]
حيث \(M(x)\) هو الحد الرئيسي و\(E(x)\) حد الخطأ. ولا يكفي ذكر أن \(E(x)\) ``صغير''؛ بل يجب تحديد معنى الصغر، ومجال الانتظام، واعتماد الثوابت على المعلمات.

\section{رمز لاندو الكبير}

\begin{definition}
لتكن \(f\) و\(g\) دالتين معرفتين عندما \(x\geq x_0\)، ولنفترض أن \(g(x)>0\) هناك. نكتب
\[
f(x)=O(g(x))
\qquad (x\to\infty)
\]
إذا وُجد ثابت \(C>0\) وعدد \(x_1\geq x_0\) بحيث
\[
|f(x)|\leq Cg(x)
\]
لكل \(x\geq x_1\).
\end{definition}

ويقال أيضًا إن
\[
f(x)\ll g(x).
\]
والرمزان \(O\) و\(\ll\) متكافئان في هذا الكتاب، ما لم نصرح بخلاف ذلك.

\begin{example}
لدينا
\[
\log x=O(x^\varepsilon)
\]
لكل \(\varepsilon>0\). لكن الثابت الضمني يعتمد عادة على \(\varepsilon\).
\end{example}

\begin{example}
من
\[
3x^2-7x+1=O(x^2)
\]
لا نستنتج أن المقدار قريب نسبيًا من \(3x^2\)، بل نستنتج فقط أنه لا يتجاوز رتبة \(x^2\).
\end{example}

\section{الاعتماد على المعلمات والانتظام}

إذا كان الثابت الضمني يعتمد على معلمة \(\alpha\)، نكتب
\[
f(x,\alpha)=O_\alpha(g(x,\alpha))
\]
أو
\[
f(x,\alpha)\ll_\alpha g(x,\alpha).
\]

\begin{definition}
نقول إن التقدير
\[
f(x,\alpha)=O(g(x,\alpha))
\]
منتظم بالنسبة إلى \(\alpha\in A\) إذا أمكن اختيار ثابت واحد \(C\) صالح لجميع \(\alpha\in A\).
\end{definition}

الانتظام ليس تفصيلًا شكليًا. ففي مسائل الأوليات في المتتاليات الحسابية، قد يكون التقدير صحيحًا لكل مقياس \(q\) ثابت، لكنه عديم الفائدة إذا ازداد الثابت بسرعة عندما يكبر \(q\).

\section{رمز لاندو الصغير}

\begin{definition}
نكتب
\[
f(x)=o(g(x))
\qquad (x\to\infty)
\]
إذا
\[
\lim_{x\to\infty}\frac{f(x)}{g(x)}=0.
\]
\end{definition}

إذن \(o(g)\) أقوى من \(O(g)\). فإذا كان \(f=o(g)\)، فإن \(f=O(g)\)، لكن العكس غير صحيح عمومًا.

\begin{example}
لدينا
\[
\log x=o(x^\varepsilon)
\]
لكل \(\varepsilon>0\)، و
\[
x=o(x\log x),
\]
لكن
\[
x\neq o(x).
\]
\end{example}

\section{التكافؤ التقاربي}

\begin{definition}
نكتب
\[
f(x)\sim g(x)
\qquad (x\to\infty)
\]
إذا
\[
\lim_{x\to\infty}\frac{f(x)}{g(x)}=1.
\]
\end{definition}

وهذا يكافئ
\[
f(x)=g(x)+o(g(x)).
\]

\begin{proposition}
إذا كان \(f(x)\sim g(x)\) و\(g(x)\sim h(x)\)، فإن
\[
f(x)\sim h(x).
\]
\end{proposition}

\begin{proof}
لدينا
\[
\frac{f(x)}{h(x)}
=
\frac{f(x)}{g(x)}
\frac{g(x)}{h(x)},
\]
وكل عامل في الطرف الأيمن يؤول إلى \(1\).
\end{proof}

\begin{remark}
التكافؤ التقاربي لا يحفظ دائمًا العمليات غير الخطية دون شروط إضافية. فمثلًا، من \(f\sim g\) لا يلزم تلقائيًا أن \(\log f\sim\log g\)، لأن ذلك يعتمد على رتبة نمو اللوغاريتمين.
\end{remark}

\section{رموز أخرى شائعة}

نكتب
\[
f(x)\asymp g(x)
\]
إذا كان
\[
f(x)\ll g(x)
\qquad\text{و}\qquad
g(x)\ll f(x).
\]

أما الرمز
\[
f(x)=\Omega(g(x))
\]
فيستخدم بأكثر من اصطلاح في الأدبيات. لذلك ستصرح الموسوعة دائمًا بالمعنى المقصود: هل هو نفي \(o(g)\)، أم وجود متتالية تحقق حدًا سفليًا، أم حد من جهة واحدة.

\section{قواعد حسابية لرموز الرتبة}

\begin{proposition}
إذا كان \(f_1=O(g_1)\) و\(f_2=O(g_2)\)، فإن
\[
f_1+f_2=O(|g_1|+|g_2|)
\]
و
\[
f_1f_2=O(g_1g_2).
\]
\end{proposition}

\begin{proof}
تتبع النتيجتان مباشرة من متباينة المثلث وضرب الحدين.
\end{proof}

\begin{proposition}
إذا كان \(f=o(g)\) و\(h=O(k)\)، فإن
\[
fh=o(gk).
\]
\end{proposition}

\begin{proof}
نكتب
\[
\frac{|f(x)h(x)|}{|g(x)k(x)|}
=
\frac{|f(x)|}{|g(x)|}\cdot
\frac{|h(x)|}{|k(x)|}.
\]
العامل الأول يؤول إلى الصفر، والثاني محدود.
\end{proof}

\section{مقارنة القوى واللوغاريتمات}

\begin{theorem}
لكل \(\alpha>0\) ولكل عدد حقيقي \(A\)، لدينا
\[
(\log x)^A=o(x^\alpha)
\qquad (x\to\infty).
\]
\end{theorem}

\begin{proof}
يكفي عند \(A>0\) وضع \(x=e^t\). عندئذ تصبح النسبة
\[
\frac{(\log x)^A}{x^\alpha}
=
\frac{t^A}{e^{\alpha t}},
\]
وهي تؤول إلى الصفر لأن الدالة الأسية تتغلب على كل قوة. أما إذا كان \(A\leq 0\)، فالنتيجة أبسط.
\end{proof}

هذا السلم من الرتب يظهر باستمرار:
\[
1\ll \log\log x\ll \log x\ll (\log x)^A\ll x^\varepsilon\ll x^\delta
\]
عندما تكون \(1<A\) و\(0<\varepsilon<\delta\)، مع تفسير كل علاقة في المجال المناسب.

\section{المجاميع الجزئية}

لتكن \((a_n)\) متتالية عقدية، وعرّف
\[
A(x)=\sum_{n\leq x}a_n.
\]
تسمى \(A(x)\) دالة المجاميع الجزئية للمتتالية.

الفكرة الأساسية هي أن المعلومات عن \(A(x)\) يمكن نقلها إلى مجموع موزون
\[
\sum_{n\leq x}a_nf(n)
\]
إذا كانت \(f\) دالة منتظمة بما يكفي.

\section{صيغة الجمع الجزئي المنفصلة}

\begin{theorem}[الجمع الجزئي]
لتكن \(a_n\) و\(b_n\) متتاليتين، ولتكن
\[
A_n=\sum_{k=1}^{n}a_k.
\]
عندئذ
\[
\sum_{n=1}^{N}a_nb_n
=
A_Nb_N
+
\sum_{n=1}^{N-1}A_n(b_n-b_{n+1}).
\]
\end{theorem}

\begin{proof}
بما أن \(a_n=A_n-A_{n-1}\)، مع \(A_0=0\)، فإن
\[
\sum_{n=1}^{N}a_nb_n
=
\sum_{n=1}^{N}(A_n-A_{n-1})b_n.
\]
نفصل المجموعين ثم نغير الفهرس في الثاني:
\[
\sum_{n=1}^{N}A_nb_n
-
\sum_{n=0}^{N-1}A_nb_{n+1}.
\]
حد \(n=0\) يساوي صفرًا، وبجمع الحدود المشتركة نحصل على
\[
A_Nb_N+\sum_{n=1}^{N-1}A_n(b_n-b_{n+1}).
\]
\end{proof}

هذه الصيغة هي النظير المنفصل للتكامل بالتجزئة.

\section{صيغة أبيل المستمرة}

\begin{theorem}[الجمع الجزئي لأبيل]
لتكن
\[
A(t)=\sum_{n\leq t}a_n,
\]
ولتكن \(f\) دالة من الصنف \(C^1\) على \([1,x]\). عندئذ
\[
\sum_{n\leq x}a_nf(n)
=
A(x)f(x)
-
\int_1^x A(t)f'(t)\,dt.
\]
\end{theorem}

\begin{proof}
الدالة \(A(t)\) درجية، وتقفز بمقدار \(a_n\) عند كل عدد صحيح \(n\). لذلك يمكن كتابة المجموع كتکامل ستيلتج:
\[
\sum_{n\leq x}a_nf(n)
=
\int_{1^-}^{x}f(t)\,dA(t).
\]
وبالتكامل بالتجزئة لستيلتج:
\[
\int_{1^-}^{x}f\,dA
=
A(x)f(x)-\int_1^x A(t)f'(t)\,dt.
\]
ويمكن أيضًا اشتقاق الصيغة مباشرة من النسخة المنفصلة بتطبيق مبرهنة القيمة الأساسية للتفاضل على الفروق \(f(n)-f(n+1)\).
\end{proof}

\section{مثال: المجموع التوافقي}

خذ \(a_n=1\)، فيكون
\[
A(t)=\lfloor t\rfloor.
\]
وباختيار \(f(t)=1/t\)، نحصل على
\[
\sum_{n\leq x}\frac1n
=
\frac{\lfloor x\rfloor}{x}
+
\int_1^x\frac{\lfloor t\rfloor}{t^2}\,dt.
\]
وبما أن
\[
\lfloor t\rfloor=t+O(1),
\]
ينتج
\[
\sum_{n\leq x}\frac1n
=
\log x+O(1).
\]

وبتحليل أدق يمكن إثبات
\[
\sum_{n\leq x}\frac1n
=
\log x+\gamma+O\left(\frac1x\right),
\]
حيث \(\gamma\) ثابت أويلر.

\section{مثال: تحويل معلومة عدّ إلى مجموع موزون}

افترض أن
\[
A(x)=cx+O(x^\theta)
\]
لثابت \(c\) و\(0\leq\theta<1\). نريد تقدير
\[
\sum_{n\leq x}\frac{a_n}{n}.
\]
باختيار \(f(t)=1/t\)، نحصل على
\[
\sum_{n\leq x}\frac{a_n}{n}
=
\frac{A(x)}{x}
+
\int_1^x\frac{A(t)}{t^2}\,dt.
\]
وبالتعويض:
\[
\frac{A(x)}x=c+O(x^{\theta-1}),
\]
كما أن
\[
\int_1^x\frac{A(t)}{t^2}\,dt
=
c\log x
+
O\left(\int_1^x t^{\theta-2}\,dt\right).
\]
ولأن \(\theta<1\)، فإن التكامل الأخير محدود، ومن ثم
\[
\sum_{n\leq x}\frac{a_n}{n}
=
c\log x+O(1).
\]

هذه إحدى أهم وظائف الجمع الجزئي: تحويل معرفة مجموع غير موزون إلى معرفة مجموع ذي وزن سلس.

\section{نسخة على فترة}

إذا كان \(1\leq y<x\)، فإن
\[
\sum_{y<n\leq x}a_nf(n)
=
A(x)f(x)-A(y)f(y)
-
\int_y^x A(t)f'(t)\,dt.
\]
وتستخدم هذه النسخة بكثرة في الفترات القصيرة وفي تقسيم المجاميع إلى كتل.

\section{مبدأ إزالة الوزن}

إذا كان الوزن \(f\) أملس ويتغير ببطء، فإن الجمع الجزئي يسمح بنقل صعوبة المسألة إلى المجاميع الجزئية لـ\(a_n\). وعلى العكس، إذا كنا نعرف مجموعًا موزونًا مناسبًا، فقد نستعيد منه معلومة عدّ باستخدام اختيار ذكي للوزن أو تقريب دالة المؤشر.

لكن إزالة الوزن ليست مجانية دائمًا؛ فقد تفقد قوة في حد الخطأ، خصوصًا إذا كان مشتق الوزن كبيرًا أو إذا كانت المعلومات عن \(A(t)\) غير منتظمة.

\section{تكاملات ستيلتج في نظرية الأعداد}

تسمح الكتابة
\[
\sum_{n\leq x}a_nf(n)
=
\int_{1^-}^{x}f(t)\,dA(t)
\]
بالتعامل مع المجاميع والتكاملات ضمن لغة واحدة. ومن الأمثلة المهمة:
\[
\pi(x)=\int_{2^-}^{x}d\pi(t),
\]
و
\[
\sum_{p\leq x}\log p
=
\int_{2^-}^{x}\log t\,d\pi(t).
\]

ثم يؤدي التكامل بالتجزئة إلى علاقات بين دوال عدّ الأوليات المختلفة.

\section{مثال تمهيدي مرتبط بالأعداد الأولية}

إذا عرفنا تقديرًا لدالة
\[
\vartheta(x)=\sum_{p\leq x}\log p,
\]
فيمكن تحويله إلى تقدير لـ\(\pi(x)\). بالفعل،
\[
\pi(x)
=
\int_{2^-}^{x}\frac{1}{\log t}\,d\vartheta(t).
\]
وبالتكامل بالتجزئة:
\[
\pi(x)
=
\frac{\vartheta(x)}{\log x}
+
\int_2^x\frac{\vartheta(t)}{t(\log t)^2}\,dt.
\]
إذا كان
\[
\vartheta(x)\sim x,
\]
فإن هذه الصيغة تقود إلى
\[
\pi(x)\sim\frac{x}{\log x}.
\]
سيظهر البرهان الكامل، مع ضبط حدود الخطأ، في فصل مبرهنة الأعداد الأولية.

\section{أخطاء شائعة}

\begin{enumerate}
  \item إخفاء اعتماد الثابت الضمني على معلمة مهمة.
  \item استعمال \(f=O(g)\) كما لو كان \(f/g\) يملك نهاية.
  \item استبدال \(O(g_1)+O(g_2)\) بـ\(O(\min(g_1,g_2))\) بدلًا من الرتبة الأكبر.
  \item إجراء تكامل جزئي دون التحقق من انتظام الوزن.
  \item تطبيق تقدير غير منتظم داخل تكامل على مجال تعتمد حدوده على المعلمات.
  \item نسيان مساهمة الحد الطرفي \(A(x)f(x)\).
\end{enumerate}

\section{تمارين}

\begin{enumerate}
  \item أثبت أن \(f\sim g\) يكافئ \(f=g+o(g)\)، بشرط أن \(g\) لا ينعدم في النهاية.
  \item أثبت أن
  \[
  \log x=O_\varepsilon(x^\varepsilon)
  \]
  وحدد ثابتًا صريحًا صالحًا عندما \(x\geq1\).
  \item استعمل الجمع الجزئي لإثبات
  \[
  \sum_{n\leq x}n^\alpha
  =
  \frac{x^{\alpha+1}}{\alpha+1}+O(x^\alpha)
  \]
  عندما \(\alpha>0\).
  \item إذا كان \(A(x)=O(x^\theta)\)، أثبت أن
  \[
  \sum_{n\leq x}\frac{a_n}{n^\sigma}
  \]
  يبقى محدودًا عندما \(\sigma>\theta\).
  \item اشتق العلاقة بين \(\pi(x)\) و\(\vartheta(x)\) بدقة مع مراعاة الحد السفلي.
  \item أعط مثالًا لتقدير صحيح لكل معلمة ثابتة لكنه غير منتظم في المعلمة.
  \item أثبت النسخة المنفصلة من الجمع الجزئي ثم استخرج منها نسخة الفترة.
\end{enumerate}

\section{خلاصة الفصل}

اللغة التقاربية هي البنية النحوية التي تصاغ بها معظم نتائج نظرية الأعداد التحليلية، والجمع الجزئي لأبيل هو أحد أكثر أدواتها استعمالًا. فهو ينقل المعلومات من دوال العد إلى المجاميع الموزونة، ويربط المجاميع المنفصلة بالتكاملات، ويمهد للصيغ التي تربط \(\pi(x)\)، و\(\vartheta(x)\)، و\(\psi(x)\)، ودوال زيتا و\(L\).
